\section{The rank-five family}\label{sec:ppt3-family}

We first separate two tasks: enforce PPT by symmetry, then exclude biseparability using only product vectors in the
state's support. Appendix~\ref{sec:ppt3-app-family} supplies the coordinate proofs.

Let $\omega=e^{2\pi i/3}$. In the computational basis of $(\mathbb C^2)^{\otimes3}$ define the five
\emph{unnormalised} vectors
\begin{align}
u&=\ket{000}-\ket{011}-\ket{101}-\ket{110}, & v&=\ket{111}, & w&=\ket{001}+\ket{010}+\ket{100},\notag\\
e&=\ket{011}+\omega\ket{101}+\omega^2\ket{110}, & o&=\ket{100}+\omega\ket{010}+\omega^2\ket{001}. &&\label{eq:ppt3-vectors}
\end{align}
Here $u/2$ is a local-Clifford rotation of the GHZ state and $w,e,o$ are locally related $W$ states.
Global parity is $Z^{\otimes3}$, with $Z=\operatorname{diag}(1,-1)$; it has eigenvalue $+1$ on $u,e$ (an even number of
ones) and $-1$ on $v,w,o$. The cyclic permutation $\Pi$ of the qubits, $\Pi\ket{abc}=\ket{cab}$, fixes $u,v,w$ and
multiplies $e$ and $o$ by $\omega^2$. Together, $Z^{\otimes3}$ and $\Pi$ split $(\mathbb C^2)^{\otimes3}$ into two
two-dimensional $\Pi$-invariant sectors, $\operatorname{span}\{\ket{000},\ket{011}+\ket{101}+\ket{110}\}$ and
$\operatorname{span}\{v,w\}$, and four one-dimensional sectors spanned by $e,\bar e$ (even) and $o,\bar o$ (odd).
Exchanging two qubits swaps $e\leftrightarrow\bar e$ and $o\leftrightarrow\bar o$, and a PPT state invariant under such an
exchange is separable across a cut (Corollary~\ref{cor:ppt3-symmetrisation}). The support below therefore keeps $u$ from
the first sector, all of the second, and only one orientation $e,o$ of the one-dimensional sectors.
Put $B=[u,v,w,e,o]$, an $8\times5$ matrix with $B\dg B=\operatorname{diag}(4,1,3,3,3)$, and $R=\operatorname{ran}B$.
For $t>0$ set
\begin{equation}\label{eq:ppt3-rho}
\rho_t=\frac{t\,uu\dg+3t\,vv\dg+3\,ww\dg+3\,ee\dg+oo\dg}{7(t+3)} .
\end{equation}
For a three-qubit operator $X$ write $X^{T_c}$ for its partial transpose on qubit $c\in\{A,B,C\}$; the
\emph{rank profile} of $X$ is $(\rank X,\rank X^{T_A},\rank X^{T_B},\rank X^{T_C})$. For a vector $x$, $\|x\|$ is the
Euclidean norm; for an operator $X$, $\|X\|_1$ is the trace norm, $\|X\|_2=(\tr X\dg X)^{1/2}$ the Hilbert--Schmidt norm
and $\|X\|_\infty$ the operator norm (largest singular value).

\begin{theorem}[the rank-five fully-PPT GME family]\label{thm:ppt3-family}
\cls{proved; CF-1, CF-2, CF-3, IN-1, IN-2, IN-3, IN-4, IN-5, IN-6, IN-7}
Every $\rho_t$ has rank five. It is PPT at all three single-qubit cuts exactly
when
\begin{equation}\label{eq:ppt3-interval}
t_-\le t\le t_+,\qquad t_-=\frac{15+9\sqrt{17}}{16}\approx3.257,\qquad t_+=\frac{2(\sqrt{37}-1)}{3}\approx3.389,
\end{equation}
and every state in this interval is genuinely multipartite entangled (GME). In the interior the rank profile is
$(5,8,8,8)$, and at the endpoints it is $(5,7,7,7)$.
Rank five is minimal among fully-PPT GME three-qubit states.
\end{theorem}

For minimality, any fully-PPT three-qubit state of rank at most
four is separable across a $2\times4$ cut by the low-rank PPT
theorem~\cite{KCKL,HLVC}, and is therefore biseparable.

\emph{Idea of the PPT proof.} Cyclic invariance makes the three partial transposes unitarily equivalent, and global parity
splits $\rho_t^{T_A}$ into two $4\times4$ blocks (the partial transpose commutes with the real diagonal $Z^{\otimes3}$, whose eigenspaces
are spanned by $\ket{000},\ket{011},\ket{101},\ket{110}$ and by $\ket{001},\ket{010},\ket{100},\ket{111}$). The even-block
determinant forces $t\ge t_-$. In that range the odd-block determinant requires $t\le t_+$, and the first three leading
minors of both blocks are positive on the resulting interval. Sylvester's criterion then gives PPT. At an endpoint one block drops
to rank three, which gives the profiles. The minors are in Appendix~\ref{sec:ppt3-app-family}; genuine entanglement is proved
below by a witness on the range.

\subsection{Cut-product rays and a supported witness}\label{sec:ppt3-witness}

A biseparable decomposition of a state supported on $R$ uses pure vectors in $R$ that are products across some
cut. A Hermitian form nonnegative on all these vectors can therefore exclude biseparability on $R$.
For $s>0$ put
$u_s=\ket{000}-s(\ket{011}+\ket{101}+\ket{110})$, $B_s=[u_s,v,w,e,o]$ and $R_s=\operatorname{ran}B_s$, so that $u_1=u$,
$B_1=B$, $R_1=R$. An operator supported on $R_s$ is written $B_sXB_s\dg$ with a $5\times5$ coefficient matrix $X$;
subscripts identify the corresponding state. A pure cut-product vector $B_sq$ has
projector $B_sqq\dg B_s\dg$, so its coefficient matrix is $qq\dg$ and $\tr(H\,qq\dg)=q\dg Hq$ with no metric factor
$B_s\dg B_s$.

\begin{proposition}[supported witnesses]\label{prop:ppt3-witness}
\cls{proved analytically; CF-2, CF-3, IN-2, IN-3, IN-4, SW-1, SW-2}
Let $s>0$.
\begin{enumerate}
\item[(a)] The $A|BC$-product rays in $R_s$ are the rays of $B_sq_s(z)$, $z\in\mathbb C$, with
\begin{equation}\label{eq:ppt3-curve}
q_s(z)=\bigl(2z,\ -z^2(z^2+3s),\ z^2-s,\ -z(z^2+s),\ z^2+s\bigr)^T,
\end{equation}
together with the ray of $v$ (the point $z=\infty$). The other two cuts are given by the cyclic images.
\item[(b)] Let $H_s=\operatorname{diag}(-12s^3,\ 1,\ 11s^2,\ -6s,\ 14s^2)$ in the basis $(u_s,v,w,e,o)$. For
$z=x+iy$ and $r=|z|^2$,
\begin{equation}\label{eq:ppt3-sos}
q_s(z)\dg H_sq_s(z)=(r-s)^2\bigl[(r-5s)^2+12s\,x^2\bigr]\ \ge0 .
\end{equation}
At $s=1$ this is the witness $H_5:=H_1$ used for the first family,
\begin{equation}\label{eq:ppt3-h5}
H_5=\operatorname{diag}(-12,1,11,-6,14),
\qquad
q_1(z)\dg H_5q_1(z)=(r-1)^2\bigl[(r-5)^2+12x^2\bigr].
\end{equation}
Since $H_s$ is diagonal, it is invariant under the cyclic phases acting on $(u_s,v,w,e,o)$. Consequently
$\tr(H_sX)\ge0$ whenever $B_sXB_s\dg$ is biseparable, including mixtures across different cuts;
no symmetry hypothesis on $X$ is needed.
\end{enumerate}
\end{proposition}

\emph{Idea of the proof.} (a) A vector of $R_s$ is a product across $A|BC$ with first-qubit ray $(1,z)^T$ exactly when
its two $BC$ rows are proportional with ratio $z$: four linear equations in the five coordinates, with a one-dimensional
solution $q_s(z)$ for each finite $z$. (b) Since $H_s$ is diagonal, only the squared moduli of the entries of $q_s(z)$
enter, and the sum collapses to a perfect-square factorisation; a rescaling removes $s$. Diagonal forms commute with the
cyclic phases, so the other cuts give the same values, and every positive decomposition of an operator supported on $R_s$
uses vectors in $R_s$, so mixtures across cuts are controlled. The computations are in Appendix~\ref{sec:ppt3-app-family}.

\begin{proof}[Proof of the genuine entanglement in Theorem~\ref{thm:ppt3-family}]
Write $\rho_t=BX_tB\dg$. Then
\begin{equation}\label{eq:ppt3-h5-value}
\tr(H_5X_t)=\frac{29-9t}{7(t+3)}<0\qquad\text{whenever }t>29/9,
\end{equation}
and $t_->29/9$. Proposition~\ref{prop:ppt3-witness}(b) therefore proves GME throughout \eqref{eq:ppt3-interval}.
\end{proof}

For a concrete member, take $t=10/3$. Then the normalised versions of the orthogonal vectors
$(u,v,w,e,o)$ occur with probabilities
$(40,30,27,27,9)/133$, all leading minors in
\eqref{eq:ppt3-minors} are positive, and
$\tr(H_5X_{10/3})=-3/133$. Thus this fully-PPT GME state is certified by hand.
Appendix~\ref{sec:ppt3-kernel} lifts the supported witness to a global witness with an explicit noise radius around
this state (Remark~\ref{rem:ppt3-hand-radius}), again without a semidefinite-program solver.

\subsection{Reflection symmetry and local chirality}\label{sec:ppt3-chirality}

The states select one cyclic phase orientation of the $W$ vectors $e,o$; exchanging two qubits reverses it.
The following obstruction explains why this asymmetry matters.

\begin{corollary}[reflection symmetry destroys the construction]\label{cor:ppt3-symmetrisation}
\cls{proved; recalled swap-PPT reduction~\cite{NMG}}
\begin{enumerate}[label=(\alph*)]
\item A three-qubit state that is PPT at $A|BC$ and commutes with the swap of $B$ and $C$ is separable across $A|BC$.
Hence no fully-PPT GME state is invariant under the exchange of two qubits.
\item For every fully-PPT three-qubit state $\rho$, the permutation average
$\mathcal S(\rho)=\frac16\sum_{\pi\in S_3}U_\pi\rho U_\pi^\dagger$ is separable across each of the three cuts, so
\begin{equation}\label{eq:ppt3-asymmetry-upper}
 D(\rho,\mathrm{BS})\le\tfrac12\|\rho-\mathcal S(\rho)\|_1 ,
\end{equation}
where $D(\rho,\sigma)=\frac12\|\rho-\sigma\|_1$ is the trace distance, $D(\rho,\mathrm{BS})=\inf_{\sigma\in\mathrm{BS}}D(\rho,\sigma)$ and
$\mathrm{BS}$ is the set of biseparable states.
If $\rho$ is cyclically invariant and $S$ is any transposition of two qubits, the bound is
$\frac14\|\rho-S\rho S^\dagger\|_1$.
\item For any three-qubit $\rho$ that is PPT at $A|BC$, pinching with the triplet and singlet projectors $P_3$, $P_1$ of $BC$
gives an explicit $A|BC$-separable state, and
\begin{equation}\label{eq:ppt3-pinch}
D(\rho,\mathrm{BS})\le\|P_3\rho P_1\|_1 .
\end{equation}
\end{enumerate}
These are upper bounds on the distance to $\mathrm{BS}$, not entanglement tests. Coherence between singlet and triplet,
and the breaking of every pair-exchange symmetry, are therefore necessary for fully-PPT GME.
\end{corollary}
\begin{proof}
(a) The swap-invariant state is block diagonal in the triplet ($\dim3$) and singlet ($\dim1$) subspaces of $BC$, and its
partial transpose on $A$ preserves these blocks. Each block is PPT on $2\times3$ or $2\times1$, where PPT implies
separability; this is the mechanism of Novo, Moroder and G\"uhne~\cite[Sec.~III.2]{NMG}. (b) $\mathcal S(\rho)$ is fully
PPT and invariant under every transposition, so (a) applies at each cut; $\mathcal S(\rho)\in\mathrm{BS}$ gives
\eqref{eq:ppt3-asymmetry-upper}. For cyclic $\rho$, $\mathcal S(\rho)=(\rho+S\rho S^\dagger)/2$. (c) The pinched state
$P_3\rho P_3+P_1\rho P_1$ is PPT at $A|BC$ (the partial transpose on $A$ commutes with the pinching on $BC$) and block diagonal as
in (a), hence separable; its trace distance from $\rho$ is at most $\|P_3\rho P_1\|_1$. Since the swap is
$S_{BC}=P_3-P_1$, the pinched state equals $(\rho+S_{BC}\rho S_{BC})/2$ and $\|P_3\rho P_1\|_1=\frac14\|\rho-S_{BC}\rho S_{BC}\|_1$,
so (c) extends the cyclic bound of (b) to every state that is PPT at $A|BC$.
\end{proof}

The obstruction concerns reflection alone. The real vectors $u_s,v,w$ and the pair $e,o$ of \eqref{eq:ppt3-vectors} are invariant under the
antiunitary map $K\circ S_{BC}$, the swap of $B$ and $C$ followed by complex conjugation in the computational basis
(for example $S_{BC}e=\ket{011}+\omega^2\ket{101}+\omega\ket{110}=\bar e$). Hence every state with real coefficients in
these vectors, including $\rho_t$ and all states of Section~\ref{sec:ppt3-surface}, is invariant under this combined
reflection and conjugation. The GME states among them are invariant under neither separately: not under the reflection by
part~(a), and not under conjugation, which exchanges $e$ and $o$ with $\bar e$ and $\bar o$, outside their supports.

\begin{corollary}[intrinsic local chirality]\label{cor:ppt3-chirality}
\cls{proved; CF-4, SC-1--SC-5}
\begin{enumerate}[label=(\alph*)]
\item Throughout \eqref{eq:ppt3-interval}, $\rho_t$ and its complex conjugate are not related by local unitaries. The same
holds for physical spin time reversal, which differs from complex conjugation by a local unitary.
\item For every $s>0$, no invertible local filter $F_A\otimes F_B\otimes F_C$ maps $R_s$ onto its complex conjugate. Hence no state
with range $R_s$, in particular no $\rho_t$, $\sigma_6$, $\rho_{\tau,a}$ or state of the phase-coexistence line, is related to its complex
conjugate or its spin time reversal by an invertible local filter. The stabiliser of $R_s$ in $\mathrm{PGL}(2,\mathbb C)^{\times3}$ is
isomorphic to $S_3$, acting by the same M\"obius map on the three qubits.
\end{enumerate}
\end{corollary}
\emph{Idea of the proof.} (a) The eigenrays of $u$ and $v=\ket{111}$ are nondegenerate, so a local unitary relating $\rho_t$
to $\overline{\rho_t}$ must fix both, leaving only $I$ and $Z^{\otimes3}$; both fix $\rho_t$, which is not real
($\langle011|\rho_t|101\rangle\propto t+3\omega^2$). (b) $R_s$ contains exactly three fully product rays, $[\ket1^{\otimes3}]$ and
$[(\ket0\pm i\sqrt s\ket1)^{\otimes3}]$, and a filter must permute them; a M\"obius map is fixed by its action on three points, so the filter is
$F_0^{\otimes3}$ up to a scalar, which commutes with the cyclic shift. The shift has eigenvalue $\omega^2$ on $e,o\in R_s$ but not on $\overline{R_s}$.
See Appendix~\ref{sec:ppt3-app-family}.

\subsection{\texorpdfstring{An extreme point of profile $(5,6,6,6)$}{An extreme point of profile (5,6,6,6)}}\label{sec:ppt3-extreme6}

Adjusting the weight of $v$ makes both partial-transpose parity blocks singular and yields an extreme fully-PPT state.

\begin{corollary}[an extreme fully-PPT GME state of profile $(5,6,6,6)$]\label{cor:ppt3-extreme6}
\cls{exact-symbolic certificate, computer-assisted; the GME part is proved analytically; CF-6}
Set $a=t_-$ and $b=8(9-a)/5$ ($a\approx3.257$, $b\approx9.189$). Then
\begin{equation}\label{eq:ppt3-sigma6}
\sigma_6=\frac{a\,uu\dg+b\,vv\dg+3\,ww\dg+3\,ee\dg+oo\dg}{4a+b+21}
\end{equation}
is an extreme point of the set of fully-PPT three-qubit states and is GME. Its coefficient matrix is
$X_6=\operatorname{diag}(a,b,3,3,1)/(4a+b+21)$.
\end{corollary}
\emph{Idea of the proof.} The choice $a=t_-$ makes the even partial-transpose block singular, and that of $b$ makes the odd
block singular, each with a positive $3\times3$ leading minor, so every partial transpose has rank six. GME follows from the
witness $H_5$. By the supported-face criterion of~\cite{GIM}, extremality means that the only Hermitian $D$ supported on
$\operatorname{ran}\sigma_6$ whose partial transposes annihilate the kernels of those of $\sigma_6$ is a multiple of
$\sigma_6$; an exact rank computation on the $25$-dimensional space of coefficient matrices shows this
(Appendix~\ref{sec:ppt3-app-family}).
